\section{Experiments}
\subsection{Experimental Setup}
\paragraph{Datasets}  
We evaluate ReTAG on both static graphs and dynamic interaction graphs. 
The static datasets cover molecular and protein graphs with varying sizes and label settings, 
while the dynamic datasets consist of large-scale temporal interaction networks used for link prediction.

\textbf{Static Datasets.}
We consider four widely used benchmarks for graph-level and node-level prediction tasks:

\begin{itemize}
    \item \textbf{PROTEINS} \cite{protein_reference}: 
    A dataset of 1,113 protein graphs representing structural interaction patterns. 
    It is used for both node classification and graph classification tasks.

    \item \textbf{COX2} \cite{rossi2015networkrepository}: 
    A molecular graph dataset containing 467 chemical compounds related to COX-2 inhibitors. 
    It is used for binary graph classification.

    \item \textbf{ENZYMES} \cite{enzymes_reference}: 
    A benchmark of 600 enzyme graphs categorized into six functional classes. 
    It supports both node-level and graph-level classification.

    \item \textbf{BZR} \cite{rossi2015networkrepository}: 
    A dataset of 405 molecular graphs describing ligands for the benzodiazepine receptor. 
    It is used for binary graph classification.
\end{itemize}

\textbf{Dynamic Datasets.}
For temporal link prediction, we employ two large-scale dynamic interaction graphs, 
constructed from time-ordered interaction logs:
\begin{itemize}
    \item \textbf{KOUBEI}: 
    A nine-week user–store interaction graph containing hundreds of thousands of nodes and several million edges. 
    It is used for temporal edge prediction.

    \item \textbf{AMAZON}: 
    A product–user interaction graph derived from weekly Amazon review logs. 
    It contains more than two hundred thousand nodes and is used for dynamic link prediction.
\end{itemize}
A summary of the dataset statistics, including graph sizes, feature dimensions, and task types, is provided in Table~\ref{tab:datasets}.

\begin{table}[t]
\centering
\renewcommand{\arraystretch}{1.15}
\setlength{\tabcolsep}{5pt}
\footnotesize
\caption{Statistics of the experimental datasets.}
\label{tab:datasets}

\resizebox{\linewidth}{!}{%
\begin{tabular}{l cc cccc}
\toprule
\textbf{Statistics} & \multicolumn{2}{c}{\textbf{Dynamic Graphs}} & \multicolumn{4}{c}{\textbf{Static Graphs}} \\
\cmidrule(r){2-3} \cmidrule(l){4-7}
 & KOUBEI & AMAZON & PROTEINS & COX2 & ENZYMES & BZR \\
\midrule
Nodes / Graph      & 221,366   & 238,735   & 39.06 & 41.22 & 32.63 & 35.75 \\
Edges / Graph      & 3986609 & 876237   & 72.82 & 43.45 & 62.14 & 38.36 \\
Density            & 3.3e-4    & 6.2e-5    & 4.8e-2 & 2.6e-2 & 5.9e-2 & 3.0e-2 \\
Graphs             & 1         & 1         & 1,113 & 467   & 600   & 405 \\
Graph Classes      & --        & --        & 2     & 2     & 6     & 2 \\
Node Features      & --        & --        & 3     & 1     & 18    & 3 \\
\midrule
Task               & Edge      & Edge      & Node, Graph & Graph & Node, Graph & Graph \\
\bottomrule
\end{tabular}
} % end resizebox
\end{table}


% \paragraph{Comparison Methods}
% To verify the effectiveness of ReTAG, we conduct comprehensive experiments across various graph learning tasks. For dynamic graphs, we include comparisons with 
% \textit{GraphCL-based Methods}: LightGCN~\cite{he2020lightgcn}, SGL~\cite{wu2021sgl}, MixGCF~\cite{mixgcf2021}, SimGCL~\cite{simgcl2021}, 
% \textit{Prompt-based}: GraphPro~\cite{liu2023graphprompt}, GraphPro+PRODIGY~\cite{feng2023prodigy}, and 
% \textit{Retrieval-Augmented}: RAGraph~\cite{JiangQXZZ00X0W24}. 
% For static graphs, we compare ReTAG with representative models, including 
% \textit{GNN Baselines}: GCN~\cite{kipf2017gcn}, GraphSAGE~\cite{hamilton2017inductive}, GAT~\cite{velivckovic2018gat}, GIN~\cite{xu2019gin}, 
% \textit{Retrieval-Augmented Variants}: RAGraph~\cite{JiangQXZZ00X0W24}, 
% \textit{Prompt-based Variants}: GraphPrompt~\cite{xie2022graphprompt}, GraphPrompt+PRODIGY~\cite{feng2023prodigy}, and ProNoG~\cite{Yu2025ProNoG}.

\paragraph{Comparison Methods}
We compare ReTAG with representative baselines for both dynamic and static graph learning tasks.

\textbf{Dynamic Graphs.}
\begin{itemize}
    \item \textit{GraphCL-based Methods}:
    \textbf{LightGCN}~\cite{lightgcn_reference} linearly propagates embeddings for user–item recommendation. 
    \textbf{SGL}~\cite{sgl_reference} adds self-supervised discrimination to enhance representation learning. 
    \textbf{MixGCF}~\cite{mixgcf_reference} mixes synthetic negatives across layers for collaborative filtering. 
    \textbf{SimGCL}~\cite{simgcl_reference} applies contrastive learning for robust dynamic link prediction.
    \textbf{StabCF}~\cite{wu2026stabcf} proposes a stabilized training strategy for collaborative filtering that mitigates gradient oscillation in GNN-based recommenders.

    \item \textit{Prompt-based Methods}:
    \textbf{GraphPro}~\cite{graphpro_reference} introduces spatial–temporal prompts for dynamic graphs. 
    \textbf{GraphPro+PRODIGY}~\cite{prodigy_reference} enhances prompting with in-context learning.

    \item \textit{Retrieval-Augmented Methods}:
    \textbf{RAGraph}~\cite{JiangQXZZ00X0W24} retrieves similar graphs via multi-dimensional similarity 
    and injects external structural knowledge into the query graph.
\end{itemize}

\textbf{Static Graphs.}
\begin{itemize}
    \item \textit{GNN Baselines}:
    \textbf{GCN}~\cite{gcn_reference} aggregates neighborhood information for node/graph classification. 
    \textbf{GraphSAGE}~\cite{graphsage_reference} generates inductive embeddings via neighbor sampling. 
    \textbf{GAT}~\cite{gat_reference} uses attention-weighted neighbor aggregation. 
    \textbf{GIN}~\cite{gin_reference} applies MLP-based injective aggregation for expressive graph representations.

    \item \textit{Retrieval-Augmented Variant}:
    \textbf{RAGraph}~\cite{JiangQXZZ00X0W24} integrates retrieved structural priors to enhance GNN learning.

    \item \textit{Prompt-based Variants}:
    \textbf{GraphPrompt}~\cite{graphprompt_reference} employs task-specific prompts for unified graph pre-training. 
    \textbf{GraphPrompt+PRODIGY}~\cite{prodigy_reference} introduces contextual prompting with in-context examples. 
    \textbf{ProNoG}~\cite{yu2024nonhomophilic} generates node-specific prompts via multi-hop non-homophily patterns.
    \textbf{GP2F}~\cite{he2026gp2f} performs cross-domain graph prompting with adaptive fusion of multiple pre-trained GNNs for few-shot graph adaptation.
\end{itemize}





\paragraph{Evaluation Metrics}
For node and graph classification tasks, we report the classification accuracy.
For dynamic link prediction, we follow prior work and use top-$K$ ranking metrics,
including Recall@20 and nDCG@20, computed on the held-out test snapshots.
\paragraph{Parameters Settings}
We adopt a training–resource split: model training is performed on the labeled
data, while retrieval is conducted from a disjoint resource set.
For static graphs, we use a 50\%/30\%/20\%
split of nodes (train/validation/test).
For dynamic graphs, interaction snapshots are chronologically partitioned, and
only past snapshots are used as the resource set for time-aware retrieval.
Retrieval-based methods (e.g., PRODIGY, RAGraph) retrieve exclusively from
the resource set to prevent information leakage, whereas other baselines are
fine-tuned on the merged graph.
All experiments are implemented in PyTorch and run on a single NVIDIA A100
GPU (40GB).
We use the Adam optimizer with learning rate $1\times10^{-3}$ and batch size 16.
For classification tasks, the hidden dimension is set to 256, and the first stage
of MTMP employs $L=2$ layers.
Cell complexes are constructed from $k$-hop neighborhoods, where $k$ is
selected on the validation set (typically $k\in\{2,3\}$).
The CTCL module applies a 2-dimensional contrastive loss with weight
$\lambda\in[0.1,0.4]$.
For link prediction, the first stage of MTMP adopts a 3-layer architecture
($L=3$) with 64-dimensional embeddings, and the whole model is trained
end-to-end with both contrastive and retrieval modules.

\subsection{Cellular Complex-Augmented Graph Results}
% \paragraph{Overall Performance} 
The experimental results of our model and baselines on node and graph classification for static graphs and on link prediction for dynamic graphs are summarized in Table~\ref{tab:static-results} and Table~\ref{tab:dynamic-results}. Based on the comparison with existing baselines, we draw the following observations:
% \vspace{-8pt}
\setlength{\leftmargini}{0.5em}
\begin{itemize}
   \item \textbf{Our proposed ReTAG achieves the best performance.} 
ReTAG significantly outperforms all baselines on both static and dynamic graphs. 
On static datasets, ReTAG achieves a relative improvement of 2.3\%–3.7\% in node classification accuracy and 1.9\%–4.4\% in graph classification accuracy over state-of-the-art methods, consistently achieving the highest performance across all tasks. 
On dynamic benchmarks, ReTAG achieves relative gains of 2.7\%–6.7\% in Recall@20 and 2.3\%-6.2\% in nDCG@20 over strong baselines such as RAGraph/FT.
These gains can be attributed to ReTAG's ability to incorporate high-dimensional topological knowledge via cell complex lifting and retrieve structured support from external graphs. 
Its MTMP module enables hierarchical propagation across complex structures, while the CTCL loss regularizes representations by enforcing 2-cell consistency, jointly enhancing expressiveness and generalizability across tasks.



% \vspace{-2pt}
    \item \textbf{Retrieval-augmented topology reasoning is crucial for generalized graph learning.} 
From Table~\ref{tab:static-results} and Table~\ref{tab:dynamic-results}, we observe that retrieval-augmented baselines (e.g., RAGraph/FT) consistently outperform standard GNNs and prompt-based variants, demonstrating that incorporating external knowledge via retrieval is vital for generalization. 
However, these baselines often rely on shallow node- or edge-level similarity, limiting their ability to capture complex structural semantics. 
ReTAG addresses this limitation by lifting graphs into higher-dimensional cell complexes and conducting retrieval based on both semantic and topological alignment. 
This enables the model to obtain more informative and transferable support subgraphs. 
With MTMP injection and CTCL regularization, ReTAG better integrates retrieved knowledge, yielding superior performance on various graph tasks.



\end{itemize}
\begin{table*}[ht]
\centering
\renewcommand{\arraystretch}{1.15}
\caption{Accuracy evaluation on node and graph classification. All tabular results (\%) are reported as mean±std across five runs. 
The best results are \textbf{bolded} and the runner-ups are \underline{underlined}.}
\resizebox{\textwidth}{!}{
\begin{tabular}{lcccccc}
\toprule

\multirow{2}{*}{\textbf{Methods}} 
& \multicolumn{2}{c}{\textbf{Node Classification}} 
& \multicolumn{4}{c}{\textbf{Graph Classification}} \\
\cmidrule(lr){2-3}  % Node Classification 下划线
\cmidrule(lr){4-7}  % Graph Classification 下划线
& \textbf{PROTEINS (5-shot)} & \textbf{ENZYMES (5-shot)} 
& \textbf{PROTEINS (5-shot)} & \textbf{COX2 (5-shot)} 
 & \textbf{BZR (5-shot)} & \textbf{ENZYMES (5-shot)}\\
\cmidrule(lr){1-7} 
{\textit{GNN Baselines}} \\
GCN        & 46.63$\pm$03.04 & 52.80$\pm$12.89 & 54.80$\pm$06.64  & 67.87$\pm$03.39  & 58.76$\pm$05.08 & 22.67$\pm$05.20 \\
GraphSAGE  & 48.87$\pm$02.64 & 48.75$\pm$01.59  & 54.68$\pm$09.34 & 67.02$\pm$05.42  & 58.27$\pm$04.79 & 21.17$\pm$05.49 \\
GAT        & 48.13$\pm$07.90 & 47.75$\pm$01.23  & 55.82$\pm$07.31  & 64.89$\pm$03.23 & 57.04$\pm$06.70 & 20.67$\pm$03.27\\
GIN        & 49.61$\pm$01.58 & 48.82$\pm$01.58  & 56.17$\pm$08.58  & 62.77$\pm$02.85  & 56.54$\pm$04.20 & 21.10$\pm$02.53\\
\midrule
\multicolumn{7}{l}{\textit{GraphPrompt+ Variants}} \\
Vanilla/NF    & 44.88$\pm$13.17 & 48.81$\pm$01.88  & 56.68$\pm$03.63  & 53.04$\pm$04.13   & 68.77$\pm$03.44 & 36.50$\pm$03.31\\
Vanilla/FT    & 48.99$\pm$01.88  & 51.99$\pm$01.36  & 57.04$\pm$03.88  & 64.04$\pm$08.20   & 69.01$\pm$02.21 & 40.00$\pm$04.36\\
PRODIGY/NF    & 47.32$\pm$08.12  & 43.80$\pm$14.03 & 53.48$\pm$06.72  & 53.97$\pm$10.34  & 67.18$\pm$08.93 & 22.12$\pm$13.84\\
PRODIGY/FT    & 53.26$\pm$06.42  & 57.98$\pm$12.37 & 57.14$\pm$10.34 & 65.31$\pm$04.28   & 68.08$\pm$06.68 & 25.94$\pm$05.12\\
ProNoG    & 52.89$\pm$11.76  & \underline{76.48}$\pm$06.23 & 61.63$\pm$08.01 & 58.03$\pm$14.28   & 62.26$\pm$12.27 & 37.91$\pm$05.64\\
GP2F    & 58.43$\pm$02.07  & 49.69$\pm$01.35 & 58.88$\pm$01.39 & 56.73$\pm$07.68   & 58.75$\pm$10.93 & 23.25$\pm$02.78\\
\midrule
\multicolumn{7}{l}{\textit{Retrieval-Augmented Variants}} \\
RAGRAPH/NF  & 56.12$\pm$04.11  & 75.92$\pm$01.72  & 58.48$\pm$03.93  & 55.32$\pm$04.15   & \underline{77.53}$\pm$05.26 & 38.17$\pm$03.39\\
RAGRAPH/FT & \underline{58.74}$\pm$00.87  & 75.74$\pm$01.92  & \underline{62.33}$\pm$02.52  & \underline{76.60}$\pm$02.30  & 76.79$\pm$05.02 & \underline{47.71}$\pm$06.88\\
\midrule

\textbf{ReTAG(Ours)} & \textbf{60.91}$\pm$01.79  & \textbf{78.26}$\pm$01.57  & \textbf{64.30}$\pm$02.21  & 
\textbf{78.09}$\pm$02.57  & \textbf{79.01}$\pm$03.49  & \textbf{49.83}$\pm$04.72\\
\bottomrule
\end{tabular}
}

\label{tab:static-results}
\end{table*}
\begin{table}[t]
\centering
\caption{Performance Evaluation (\%) on Link Prediction.}
\label{tab:dynamic-results}

\footnotesize
\setlength{\tabcolsep}{3.2pt}
\renewcommand{\arraystretch}{0.92}

\begin{tabular}{lcccc}
\toprule
\multirow{2}{*}{\textbf{Methods}}  
 & \multicolumn{2}{c}{\textbf{KOUBEI}} 
 & \multicolumn{2}{c}{\textbf{AMAZON}} \\
\cmidrule(lr){2-3} \cmidrule(lr){4-5}
 & \textbf{Recall@20} & \textbf{nDCG@20} 
 & \textbf{Recall@20} & \textbf{nDCG@20} \\
\midrule

\multicolumn{5}{l}{\textit{GraphCL}} \\
LightGCN  & 30.21$\pm$06.45 & 22.24$\pm$05.83 & 15.07$\pm$06.48 & 06.53$\pm$02.66 \\
SGL       & 32.61$\pm$04.27 & 22.36$\pm$04.82 & 15.78$\pm$07.12 & 07.90$\pm$02.49 \\
MixGCF    & 32.06$\pm$04.20 & 22.49$\pm$06.91 & 15.24$\pm$08.98 & 07.40$\pm$03.44 \\
SimGCL    & 33.07$\pm$05.28 & 23.08$\pm$05.55 & 16.10$\pm$07.91 & 07.58$\pm$03.51 \\
StableCF    & 32.42$\pm$02.01 & 17.55$\pm$01.37 & 13.59$\pm$05.28 & 05.27$\pm$04.64 \\
\addlinespace[0.3ex]
\midrule

\multicolumn{5}{l}{\textit{GraphPro+}} \\
Vanilla/NF   & 21.31$\pm$04.59 & 15.31$\pm$03.11 & 12.56$\pm$07.45 & 06.31$\pm$03.92 \\
Vanilla/FT   & 33.96$\pm$04.13 & 24.66$\pm$02.78 & 18.14$\pm$07.55 & 08.73$\pm$03.74 \\
PRODIGY/NF   & 21.66$\pm$03.21 & 14.82$\pm$03.92 & 11.88$\pm$02.61 & 05.84$\pm$01.84 \\
PRODIGY/FT   & 33.46$\pm$04.70 & 23.28$\pm$03.40 & 16.72$\pm$04.28 & 08.09$\pm$02.66 \\
\addlinespace[0.3ex]
\midrule

\multicolumn{5}{l}{\textit{Retrieval-Augmented}} \\
RAGRAPH/NF   & 22.86$\pm$03.44 & 16.68$\pm$02.48 & 13.78$\pm$05.54 & 06.52$\pm$02.69 \\
RAGRAPH/FT   & \underline{34.27}$\pm$03.93 & \underline{24.82}$\pm$02.69 
              & \underline{18.32}$\pm$07.45 & \underline{09.09}$\pm$03.89 \\
\addlinespace[0.3ex]
\midrule

\textbf{ReTAG (Ours)}
 & \textbf{35.21}$\pm$03.21 & \textbf{25.39}$\pm$02.23
 & \textbf{19.54}$\pm$08.04 & \textbf{09.65}$\pm$04.06 \\
\bottomrule
\end{tabular}
\end{table}

















% \begin{figure}[t]
% \centering

% % ===== 左图：k-hop =====
% \begin{minipage}{0.48\linewidth}
%     \centering
%     \begin{tikzpicture}
%     \begin{axis}[
%         width=\linewidth,
%         height=4.8cm,
%         xlabel={k-hop},
%         ylabel={Accuracy (\%)},
%         ymin=60, ymax=80,
%         xtick={1,2,3,4},
%         ytick={60,62.5,65,67.5,70,72.5,75,77.5,80},
%         tick label style={font=\scriptsize},
%         label style={font=\scriptsize},
%         grid=both,
%         grid style={line width=0.4pt, dashed, gray!50},
%         axis lines=box,
%         axis line style={black, line width=0.3pt},
%         tick style={gray!60},
%         line width=0.8pt
%     ]
%     % PROTEINS
%     \addplot[blue, mark=*, thick] coordinates {
%         (1, 62.8) (2, 64.0) (3, 64.3) (4, 62.5)
%     };
%     % BZR
%     \addplot[teal, mark=triangle*, thick] coordinates {
%         (1, 78.2) (2, 78.8) (3, 79.0) (4, 78.3)
%     };
%     % COX2
%     \addplot[orange, mark=square*, thick] coordinates {
%         (1, 75.1) (2, 78.1) (3, 77.0) (4, 74.9)
%     };
%     \end{axis}
%     \end{tikzpicture}
% \end{minipage}
% \hfill
% % ===== 右图：Contrastive Loss Weight =====
% \begin{minipage}{0.48\linewidth}
%     \centering
%     \begin{tikzpicture}
%     \begin{axis}[
%         width=\linewidth,
%         height=4.8cm,
%         xlabel={CTCL Loss Weight($\alpha$)},
%         ylabel={Accuracy (\%)},
%         ymin=60, ymax=80,
%         xtick={0,0.1,0.2,0.3,0.4,0.5},
%         ytick={60,62.5,65,67.5,70,72.5,75,77.5,80},
%         legend to name=sharedlegend,
%         legend columns=3,
%         legend style={font=\scriptsize, /tikz/every even column/.append style={column sep=0.4cm}},
%         tick label style={font=\scriptsize},
%         label style={font=\scriptsize},
%         grid=both,
%         grid style={line width=0.4pt, dashed, gray!50},
%         axis lines=box,
%         axis line style={black, line width=0.3pt},
%         tick style={gray!60},
%         line width=0.8pt
%     ]
%     % PROTEINS
%     \addplot[blue, mark=*, thick] coordinates {
%         (0, 61.5) (0.1, 63.2) (0.2, 63.5) (0.3, 63.8) (0.4, 64.3) (0.5, 62.8)
%     };
%     \addlegendentry{PROTEINS}
%     % BZR
%     \addplot[teal, mark=triangle*, thick] coordinates {
%         (0, 78.2) (0.1, 78.8) (0.2, 79.0) (0.3, 78.5) (0.4, 78.2) (0.5, 77.6)
%     };
%     \addlegendentry{BZR}
%     % COX2
%     \addplot[orange, mark=square*, thick] coordinates {
%         (0, 75.5) (0.1, 76.6) (0.2, 78.1) (0.3, 77.0) (0.4, 75.7) (0.5, 75.1)
%     };
%     \addlegendentry{COX2}
%     \end{axis}
%     \end{tikzpicture}
% \end{minipage}

% \ref{sharedlegend} 

% \caption{
% Accuracy with varying k-hop size (left) and CTCL loss weight (right) on PROTEINS, BZR, and COX2.
% }
% \label{fig:hyperparam-study}
% \end{figure}

\begin{figure}[t]
    \centering
    % 子图 (a)
    \subfloat[]{
        \includegraphics[width=0.47\linewidth]{Pics/Hyper_khop.pdf}
        \label{fig:khop}
    }
    \hfil
    % 子图 (b)
    \subfloat[]{
        \includegraphics[width=0.47\linewidth]{Pics/Hyper_ctcl.pdf}
        \label{fig:ctcl}
    }

    \caption{Accuracy under different settings. 
    (a) Varying k-hop size. 
    (b) Varying CTCL weight.}
    \label{fig:analysis}
\end{figure}



% ($\alpha$)



% \vspace{-5pt}
\subsection{Hyper-parameter Sensitivity Study} 
We investigate the sensitivity of ReTAG to two key hyper-parameters: the
$k$-hop neighborhood size used for constructing cell complexes, and the
cellular topological contrastive learning (CTCL) loss weight. These
hyper-parameters respectively determine the receptive field of the lifted
complex and the strength of 2-cell topological regularization. We vary
$k$ from 1 to 4 and the CTCL loss weight from 0 to 0.5 to analyze their
individual effects. The results are presented in Fig.~\ref{fig:analysis}.

\begin{itemize}
    \item \textbf{$k$-hop Neighborhood Size.}  
    Fig.~\ref{fig:analysis}(a) shows the impact of varying the $k$-hop
    receptive field. As $k$ increases, each cell complex incorporates
    richer structural context and enables retrieval of larger supporting
    subgraphs. However, excessively large $k$ introduces noisy or redundant
    structures that degrade performance, while overly small $k$ fails to
    capture essential topological motifs for retrieval and reasoning. This
    highlights the importance of selecting a moderate $k$ to balance
    completeness and noise.

    \item \textbf{CTCL Loss Weight $\lambda$.}  
    Fig.~\ref{fig:analysis}(b) examines the effect of the CTCL loss weight.
    When $\lambda$ is too small, the model underutilizes discriminative
    2-cell information, weakening topological regularization. Conversely,
    large values of $\lambda$ overly constrain the optimization and conflict
    with the primary learning objective, resulting in reduced accuracy. A
    moderate CTCL weight therefore provides the best trade-off between
    topological consistency and task-specific learning.
\end{itemize}
% We study the sensitivity of ReTAG to two key hyper-parameters: the number of hops ($k$) used to construct cell complexes, and the cellular topological contrastive learning (CTCL) loss weight. The $k$-hop value determines the receptive field of each complex, affecting the amount of structural context for retrieval and reasoning. The contrastive loss weight controls the strength of 2-cell topological regularization during training. We vary $k$ from 1 to 4, and the loss weight from 0 to 0.5 to assess their respective impacts. 
% Fig.~\ref{fig:analysis}(a) shows the effect of different $k$-hop values on classification accuracy. As $k$ increases, the volume of retrieved knowledge expands accordingly. While larger $k$ introduces noisy or redundant substructures that hinder performance, excessively small $k$ may miss crucial topological patterns needed for effective retrieval.
%  Fig.~\ref{fig:analysis}(b) illustrates the impact of varying the CTCL loss weight. Insufficient weight underutilizes the discriminative 2-dimensional information, while too large a weight overly constrains learning, interfering with the main optimization objective and degrading downstream performance. These trends highlight the importance of balancing topological regularization with task-specific learning in ReTAG.








% \begin{table}[t]
% \centering
% \footnotesize
% \caption{Ablation study on COX2 and PROTEINS datasets.}
% \renewcommand{\arraystretch}{1.15}
% \begin{tabular}{lcc}
% \toprule
% \textbf{Model Variant} & \textbf{COX2 (\%)} & \textbf{PROTEINS (\%)} \\
% \midrule
% w/o FCL         & 76.05$\pm$02.89 & 62.33$\pm$02.52 \\
% w/o MTMP         & 76.21$\pm$01.97 & 62.45$\pm$03.19 \\
% w/o CTCL   & 76.94$\pm$01.83 & 63.21$\pm$04.79 \\
% \textbf{Full Model (ReTAG)}      & \textbf{78.09$\pm$02.57} & \textbf{64.30$\pm$02.21} \\
% \bottomrule
% \end{tabular}

% \label{tab:ablation}
% \end{table}
% \vspace{5pt}  % 可以控制表格的位置，调整这个值让表格上下移动


\subsection{Few-shot Generalization Study}
We evaluate the few-shot performance of ReTAG under 1–5 shot settings 
across four datasets (PROTEINS, BZR, ENZYMES, COX2), together with 
representative baseline methods (RAGraph, Vanilla, GAT, GIN, GCN, 
GraphSAGE). As shown in Fig.~\ref{fig:fewshot-study}, accuracy 
generally improves as more labeled samples become available, following 
the typical trend of metric-based adaptation under low-data regimes. 
Across all datasets and shot levels, ReTAG maintains consistently strong 
performance and remains competitive as supervision increases. These 
results indicate that the proposed retrieval-augmented topological 
modeling is effective in scenarios with limited labeled data.








\begin{figure}[t]
    \centering

    % 第一行 (a) (b)
    \subfloat[]{%
        \includegraphics[width=0.49\linewidth]{Pics/Few_Shot_Performance_PROTEINS.pdf}%
        \label{fig:fewshot-proteins}
    }
    \hfil
    \hspace{-0.01\linewidth}%
    \subfloat[]{%
        \includegraphics[width=0.49\linewidth]{Pics/Few_Shot_Performance_BZR.pdf}%
        \label{fig:fewshot-bzr}
    }

    \vspace{0.15cm}

    % 第二行 (c) (d)
    \subfloat[]{%
        \includegraphics[width=0.49\linewidth]{Pics/Few_Shot_Performance_ENZYMES.pdf}%
        \label{fig:fewshot-enzymes}
    }
    \hfil
    \hspace{-0.01\linewidth}%
    \subfloat[]{%
        \includegraphics[width=0.49\linewidth]{Pics/Few_Shot_Performance_COX2.pdf}%
        \label{fig:fewshot-cox2}
    }

    \caption{
Few-shot accuracy comparison of ReTAG and baseline methods 
(RAGraph, Vanilla, GAT, GIN, GCN, GraphSAGE) across four datasets:
(a) PROTEINS, (b) BZR, (c) ENZYMES, and (d) COX2. 
Results are reported under the 1–5 shot setting.
}

    \label{fig:fewshot-study}
\end{figure}




% \vspace{-9pt}
\subsection{Ablation Study}  
We evaluate three ablated variants of ReTAG: removing the Fundamental Cycle-Guided Complex Lifting (FCL), which eliminates 2-cell construction and reduces the model to an edge-based graph; discarding the Multi-dimensional Topological Message Passing (MTMP), which disables hierarchical propagation across cell dimensions; and removing the Cellular Topological Contrastive Learning (CTCL), which drops contrastive regularization over 2-cell representations. As shown in Table~\ref{tab:ablation}, we observe that:

% \vspace{-8pt}
\begin{itemize}
   \item The performance drops when removing FCL, as eliminating 2-cells via fundamental cycles reduces the model to a purely edge-based graph. This shows that high-dimensional cycles are crucial for capturing richer structural context and stronger structural reasoning.
    \item The model’s performance suffers when MTMP is discarded. Disabling MTMP restricts interactions to nodes and edges, preventing the model from reasoning over cycles and their associated higher-dimensional cells. This shows that MTMP is crucial for capturing complex topological dependencies and multi-dimensional context. %Without hierarchical propagation, information exchange is confined to the 1-skeleton, preventing nodes and edges from exploiting higher-dimensional cycles. This shows that MTMP is essential for integrating multi-dimensional context and capturing topological dependencies.
    \item The performance of the model drops when removing CTCL. Without contrastive regularization over 2-cell representations, the model fails to enforce consistency among high-dimensional structures, leading to less discriminative embeddings. This shows that CTCL is essential for improving robustness and generalization in capturing cellular topology.

\end{itemize}
\begin{table}[t]
\centering
\caption{Ablation Study on COX2 and PROTEINS Datasets.}
\label{tab:ablation}

\footnotesize
\setlength{\tabcolsep}{4pt}
\renewcommand{\arraystretch}{0.92}

\begin{tabular}{lcc}
\toprule
\textbf{Model Variant} & \textbf{COX2 (\%)} & \textbf{PROTEINS (\%)} \\
\midrule
w/o FCL   & 76.05$\pm$02.89 & 62.33$\pm$02.52 \\
w/o MTMP  & 76.21$\pm$01.97 & 62.45$\pm$03.19 \\
w/o CTCL  & 76.94$\pm$01.83 & 63.21$\pm$04.79 \\

\midrule
\textbf{Full Model (ReTAG)} 
          & \textbf{78.09}$\pm$02.57 
          & \textbf{64.30}$\pm$02.21 \\
\bottomrule
\end{tabular}
\end{table}
% w/o Retrieval  & 76.94$\pm$01.83 & 63.21$\pm$04.79 \\




\subsection{Cycle Structure and Representation Analysis}
\paragraph{Cycle Structure Analysis}
To examine when ReTAG is expected to help, we analyze the cycle statistics
of all datasets. Fig.~\ref{fig:cycle-analysis}(a) shows the fraction of
nodes that participate in at least one cycle, revealing large variations
in cycle density: molecular and protein graphs are cycle rich, whereas
some benchmarks are almost tree-like, and large recommendation graphs lie
in between with moderate cycle coverage. 
Fig.~\ref{fig:cycle-analysis}(b) further shows that the number of
fundamental cycles spans several orders of magnitude across datasets.

Despite this diversity, ReTAG consistently outperforms baselines
(Table~\ref{tab:static-results} and Table~\ref{tab:dynamic-results}),
demonstrating its robustness across varying cycle densities. When cycles are
abundant, ReTAG lifts informative motifs to 2-cells and propagates them
through MTMP; when cycles are sparse, the model naturally reduces to an
edge-centric encoder while still retaining the advantages of topology-aware retrieval.
Together with the ablation results, these observations show that
explicitly modeling fundamental cycles improves representation quality
across heterogeneous topological distributions.

\paragraph{Representation Visualization}
To further inspect how cycle-aware lifting and MTMP shape the learned
representations, we visualize graph-level embeddings using t-SNE from both
real-world and controlled benchmarks. Fig.~\ref{fig:repr-vis} first reports
the visualization on ENZYMES. The standard GCN baseline in
Fig.~\ref{fig:repr-vis}(a) produces a largely isotropic cloud, where graphs
from different classes are highly mixed and no clear global organization is
formed. This pattern suggests that the encoder mainly relies on local
neighborhood statistics and has limited ability to organize graphs by their
higher-order topological structures.

In contrast, the ReTAG embeddings in Fig.~\ref{fig:repr-vis}(b) concentrate
along several smooth manifolds with more clearly separated branches. This
indicates that ReTAG arranges enzyme graphs according to shared
cycle-aware motifs: graphs with similar configurations of fundamental
cycles are mapped to nearby positions along the same curve, whereas graphs
with distinct cycle structures tend to occupy different segments of the
manifold. Such geometry naturally arises from lifting fundamental cycles to
2-cells and propagating information jointly over nodes, edges, and 2-cells
via MTMP. Although class labels are not perfectly separated in two
dimensions due to the structural heterogeneity within each enzyme class,
ReTAG still significantly outperforms GCN in
Table~\ref{tab:static-results}, confirming that explicit higher-dimensional
cycle modeling yields more discriminative and topology-consistent graph
representations.

We further validate this observation with the controlled t-SNE study in
Fig.~\ref{fig:motivation-tsne}, generated by the motivational experiment in
which two 3-regular graph classes share the same number of nodes, edges,
and degree statistics but differ in cycle topology: $K_{3,3}$ contains no
triangles, whereas the triangular prism contains two triangular 2-cell
motifs. Under conventional message passing, GCN and GAT produce heavily
overlapped embeddings and achieve only random-level accuracy
($50.07\%\pm3.41\%$ and $49.83\%\pm2.65\%$, respectively), while GraphSAGE
shows only partial separation ($84.90\%\pm5.21\%$). Once explicit cycle
features are injected as a ReTAG-like proxy for 2-cell information, the two
classes become cleanly separated in the embedding space and the accuracy
reaches $100.00\%\pm0.00\%$. This controlled visualization complements the
ENZYMES result by isolating the source of improvement: cellular
cycle-level information provides discriminative signals that are largely
orthogonal to node/edge-level neighborhood aggregation.

\begin{figure}[t]
    \centering
    % (a) Fraction of nodes in cycles
    \subfloat[]{
        \includegraphics[width=0.46\linewidth]{Pics/cycles_ratio.pdf}
        \label{fig:cycle-coverage}
    }
    \hfil
    \hspace{0.01\linewidth}%
    % (b) Avg #cycles/graph (log scale)
    \subfloat[]{
        \includegraphics[width=0.46\linewidth]{Pics/cycles_count.pdf}
        \label{fig:cycle-count}
    }
    \caption{
    Cycle statistics across datasets.
    (a) Fraction of nodes that belong to at least one cycle.
    (b) Average number of fundamental cycles per graph (log scale).
    }
    \label{fig:cycle-analysis}
\end{figure}








\begin{figure}[t]
  \centering
  \subfloat[]{
    \includegraphics[width=0.46\linewidth]{Pics/vis_ENZYMES_graph_emb_tsne_gcn.pdf}
  }\hfill
  \hspace{-0.94\linewidth}%
  \subfloat[]{
    \includegraphics[width=0.46\linewidth]{Pics/vis_ENZYMES_graph_emb_tsne_ReTAG8.pdf}
  }
  \caption{t-SNE visualization of graph-level embeddings on ENZYMES.
  (a) The GCN baseline. 
  (b) ReTAG.}
  \label{fig:repr-vis}
\end{figure}

\begin{figure*}[t]
  \centering
  \includegraphics[width=0.95\textwidth]{Pics/motivation_tsne.pdf}
  \caption{Controlled t-SNE visualization on the motivational benchmark.
  The two classes are 3-regular graphs with identical local degree
  statistics but different cycle structures: $K_{3,3}$ has no triangles,
  whereas the triangular prism contains triangular cycles. Conventional
  message-passing models produce overlapped or only partially separated
  embeddings, while the cycle-augmented ReTAG-like variant forms clearly
  separated clusters.}
  \label{fig:motivation-tsne}
\end{figure*}

\subsection{Case Study: Topology-Aware Retrieval on COX2}
To better illustrate how ReTAG exploits cellular topologies during
retrieval, we conduct a case study on a representative COX2 molecular
graph. The left panel of Fig.~\ref{fig:case-study-cox2} shows the query
graph, where three fundamental cycles (1--3) detected by FCL are lifted
into 2-cells that serve as the foundational high-dimensional topological
primitives for retrieval.
In both the query and retrieved graphs, the yellow node marks
the \emph{master 0-cell}, whose pivotal high-degree position serves as a
structural anchor to align local neighborhoods across the lifted complexes.

For the retrieved supports (right), we highlight matching 2-cells (dark-blue) and their branching environments (light-blue nodes, dashed edges). Although external to 2-cells, these branches contextualize motifs, distinguishing cycles with similar shapes but different surroundings. By explicitly visualizing this context, we verify the consistency of topological alignment around the master anchor.

As shown in the right panel, the top-1 support closely replicates all
three query 2-cells with highly consistent branching patterns, whereas
lower-ranked supports capture only subsets of these motifs with
increasingly divergent local contexts.
This case study demonstrates that ReTAG aligns graphs by jointly
leveraging high-dimensional cellular structures and their surrounding
branching geometry, enabling retrieval to capture multi-level
topological alignment beyond local adjacency.

\begin{figure*}[t]
    \centering
    \includegraphics[width=0.95\linewidth]{Pics/COX2_case.pdf}
    \caption{Topology-aware retrieval on a molecular graph from COX2.
    Left: query graph with three fundamental cycles (1--3) extracted by
    FCL and lifted to 2-cells (shaded).
    Right: top-$4$ retrieved graphs, where shaded regions mark matching
    2-cells; dark-blue nodes denote their 1-skeleton, and light-blue
    nodes with dashed edges show surrounding branching context.
    The yellow nodes mark the aligned master 0-cells across both the
    query and retrieved graphs.}
    \label{fig:case-study-cox2}
\end{figure*}











































\subsection{Complexity and Scalability Analysis}

\paragraph{Computational Complexity}
We analyze the asymptotic time and space complexity of ReTAG.
Let $n = |V|$ and $m = |E|$ denote the numbers of nodes and edges in a
graph, $r$ the number of 2-cells (fundamental cycles) in the lifted
complex, $d$ the hidden dimension, $L$ the number of message-passing
layers, $N_{\text{res}}$ the size of the resource pool, and $K$ the
number of retrieved supports.

The backbone encoder of ReTAG follows standard message passing with a
per-layer cost of $\mathcal{O}(m d)$. Consequently, an $L$-layer
encoder requires $\mathcal{O}(L m d)$ per forward pass, which is
consistent with standard GCN and RAGraph-style baselines.
Cell-complex lifting (spanning tree extraction and fundamental cycle
basis construction) is performed offline in $\mathcal{O}(m + r)$ time
and does not participate in backpropagation.
During training and inference, the MTMP module operates on the lifted
complex with a complexity of $\mathcal{O}((m + r) d)$ per layer.
Since complexes are constructed based on local $k$-hop neighborhoods,
$r$ is effectively bounded by a small constant. Thus, the overall MTMP
cost remains linear with respect to $m$.
Topology-aware retrieval computes similarities between a $d$-dimensional
query and all $N_{\text{res}}$ resource keys to select top-$K$
supports, costing $\mathcal{O}(N_{\text{res}} d)$, while the cellular
topological contrastive loss is $\mathcal{O}(r d)$ per batch.
Therefore, the online time complexity of ReTAG is linear in terms of
the graph size and $N_{\text{res}}$, incurring only a moderate
constant-factor overhead compared to conventional GNN baselines.

\paragraph{Space Complexity}
ReTAG stores node, edge, and 2-cell embeddings alongside the resource
pool, resulting in a memory usage of
$\mathcal{O}((n + m + r + N_{\text{res}}) d)$.
Thus, both the time and space complexity of ReTAG scale linearly with
the graph size and the number of retrieved resources.
% \begin{table}[t]
% \centering
% \caption{Theoretical complexity and empirical efficiency comparison of different models.}
% \label{tab:complexity}
% \footnotesize
% \setlength{\tabcolsep}{4pt}
% \renewcommand{\arraystretch}{0.92}
% \begin{tabular}{lcccc}
% \toprule
% \multirow{2}{*}{\textbf{Model}} 
% & \multirow{2}{*}{\textbf{Time}} 
% & \multirow{2}{*}{\textbf{Space}} 
% & \multicolumn{2}{c}{\textbf{Train Time (s/epoch)}} \\
% \cmidrule(lr){4-5}
% & & & \textbf{COX2} & \textbf{KOUBEI} \\
% \midrule
% GCN 
% & 1
% & 1 
% & -- & -- \\
% RAGraph 
% & 1
% & 1
% & -- & -- \\
% \midrule
% \textbf{ReTAG (ours)} 
% & 1
% & 1 
% & \textbf{--} & \textbf{--} \\
% \bottomrule
% \end{tabular}
% \end{table}
